% Original: PDF strana 26, gornji slajd.
\slajdid{02-s051}
\begin{frame}[label=02-s051]{Statički hazard: izlaz treba da ostane 1}
\setlength{\parskip}{0pt}
% element: 02-s051:e01
Radi uočavanja pojave pretpostavljamo da su sva kola idealno brza, osim invertora, čije je kašnjenje \(t_p\).
\par\vspace{2mm}
\begin{columns}[T,onlytextwidth]
\begin{column}{.52\textwidth}
% element: 02-s051:e02
{\centering\begin{tikzpicture}[circuit logic US,DEoznaka,x=1cm,y=1cm,
 gate/.style={draw,logic gate inputs=nn,minimum width=12mm,minimum height=9mm,scale=.6}]
\node[and gate US,gate] (u) at (2,1.1) {};
\node[and gate US,gate] (d) at (2,-.1) {};
\node[or gate US,gate] (f) at (3.8,.5) {};
\node[not gate US,draw,minimum width=8mm,minimum height=8mm,scale=.6] (i) at (.6,.55) {};
\draw[DEvod] (u.input 1)--(-.5,0 |- u.input 1) node[left] {$C$};
\draw[DEvod] (i.input)--(-.5,.55) node[left] {$B$};
\draw[DEvod] (i.output)--(1.25,.55)|-(u.input 2);
\draw[DEvod] (-.2,.55)|-(d.input 1);
\node[junction] at (-.2,.55) {};
\draw[DEvod] (d.input 2)--(-.5,0 |- d.input 2) node[left] {$A$};
\draw[DEvod] (u.output)--(2.8,1.1)|-(f.input 1);
\draw[DEvod] (d.output)--(2.8,-.1)|-(f.input 2);
\draw[DEvod] (f.output)--++(.4,0) node[right] {$F$};
\end{tikzpicture}
\par}
\par\vspace{2mm}
Početno stanje: \(C=B=A=1\), pa je \(F=C\bar B+BA=1\).
\par\vspace{1mm}
Menja se samo \(B:1\to0\). Konačno stanje: \(C=A=1\), \(B=0\), pa je opet \(F=1\).
\par\vspace{1mm}
Izlaz \textbf{treba da ostane 1}, \alert{\textbf{ali}} tokom prelaza javlja se \alert{\textbf{lažna nula}}.
\end{column}
\begin{column}{.43\textwidth}\centering
% element: 02-s051:e03
% Kompaktan vektorski dijagram; logičke vrednosti i redosled događaja su iz originala.
\begin{tikzpicture}[x=1cm,y=1cm,font=\fontsize{9}{10.5}\selectfont]
\foreach \r/\lab in {0/A,1/B,2/{\bar B},3/C,4/A,5/{BA},6/{C\bar B},7/F}{
 \pgfmathsetmacro{\yy}{-.50*\r}
 \draw[black!55,thin,-{Latex[length=.9mm]}] (0,\yy-.03)--(0,\yy+.39) node[left] {$\lab$};
 \draw[black!55,thin,-{Latex[length=.9mm]}] (-.08,\yy)--(3.72,\yy) node[below] {$t$};
}
\foreach \r in {0,3,4}{
 \pgfmathsetmacro{\yy}{-.50*\r}
 \draw[DEvod] (0,\yy+.32)--(3.35,\yy+.32);
}
\foreach \r in {1,5}{
 \pgfmathsetmacro{\yy}{-.50*\r}
 \draw[DEvod] (0,\yy+.32)--(.95,\yy+.32)--(1.06,\yy+.10)--(3.35,\yy+.10);
}
\foreach \r in {2,6}{
 \pgfmathsetmacro{\yy}{-.50*\r}
 \draw[DEvod] (0,\yy+.10)--(1.76,\yy+.10)--(1.88,\yy+.32)--(3.35,\yy+.32);
}
\draw[DEvod] (0,-3.18)--(.95,-3.18)--(1.06,-3.40)--(1.76,-3.40)--(1.88,-3.18)--(3.35,-3.18);
% Trajanje je iznad svih signala, da strelica ne seče talasne oblike.
\draw[DEvod,dashed] (1,.63)--(1,-3.57);
\draw[DEvod,dashed] (1.82,.63)--(1.82,-3.57);
\draw[DEvod,{Latex[length=1.1mm]}-{Latex[length=1.1mm]}] (1,.72)--(1.82,.72)
 node[midway,above,inner sep=1pt] {$t_p$};
\node[inner sep=1pt] at (2.7,-3.83) {lažna nula};
\end{tikzpicture}

\end{column}
\end{columns}
\end{frame}
